\section{误差分析与敏感性讨论}

\subsection{类别级别的识别差异}

从当前分类报告可以看到，不同类别在 precision、recall 与 F1-score 上并不均衡。虽然 Accuracy 能够反映整体水平，但对于分子亚型任务而言，更重要的是模型是否能够稳定识别每一个亚型。部分类别在当前划分中支持样本较少，因此一旦测试集切分后该类样本数量进一步减少，就很容易出现召回率波动较大的现象。

这类问题在生物医学分类中非常常见，因为临床数据本身就带有天然的不平衡性。某些亚型在真实世界中就较少出现，而某些亚型虽然样本较多，却可能与其他亚型存在边界重叠。因此，单次实验中某一类的表现偏低，并不一定意味着模型完全失败，而可能只是样本划分导致的统计波动。

\subsection{特征数量对结果的影响}

高维特征筛选的一个重要作用是控制模型的有效复杂度。若保留的特征数过少，模型可能无法充分表达亚型差异；若特征数过多，则模型容易受到噪声干扰，尤其在支持样本有限的情况下，训练集表现会显著好于测试集。

当前“不做特征选择”的消融实验结果就说明了这一点。虽然理论上保留全部特征可以避免信息丢失，但在实际分类任务中，这种“保守”策略往往并不优，因为大多数特征并不携带稳定判别信息。通过高方差筛选，将模型注意力集中到更有差异的基因或探针上，反而更容易得到稳健结果。

\subsection{融合策略的敏感性}

Concat 和 MOFA 的比较也反映出多组学融合方法对参数和样本量的敏感性。Concat 的优势是确定性强，只要特征维度和样本顺序一致，结果通常比较稳定。MOFA 则引入了潜在因子学习过程，其性能不仅受到输入维度影响，还受到因子数、初始化、稀疏约束与训练收敛情况影响。

因此，在当前实验规模下，Concat 可能更容易跑出较高的单次准确率，而 MOFA 的价值更多体现在可解释性与结构提取能力上。若后续扩大样本规模、细化超参数搜索，MOFA 的优势很可能更加明显。

\subsection{训练测试划分的敏感性}

训练-测试划分方式对于本研究尤其重要。由于样本总数有限，每一次随机切分都会改变类别分布和难例分布。对于少数类而言，几个样本的变化就可能影响最终指标。因此，单次实验可以作为方法验证，但不宜直接作为最终结论。

从更严谨的角度看，理想做法是使用多次重复划分、交叉验证或嵌套验证来报告统计结果。这样既可以评估平均性能，也可以估计方法对数据划分的稳定性。当前报告将这一点作为后续扩展方向，是因为目前代码结构已经具备支撑重复实验的基础，只需要进一步增加自动化统计即可。

\subsection{错误来源的分类}

对当前项目而言，错误来源大致可以分为四类：
\begin{enumerate}
  \item 数据错误：如样本名称不一致、临床标签格式不统一；
  \item 预处理错误：如特征筛选在错误阶段执行；
  \item 接口错误：如 MOFA 输入格式与第三方库预期不匹配；
  \item 评估错误：如测试集长度与预测结果长度不一致。
\end{enumerate}

本文在开发过程中实际上已经遭遇并修复了这些问题，因此将它们写入报告具有很强的实践意义。换句话说，误差分析不只是对模型预测错误的解释，也是在总结“系统性错误是如何被避免的”。

\subsection{结果稳定性的解释}

当前结果虽然还不能说具有大规模统计意义，但已经表现出两个值得注意的趋势：一是多组学融合并没有失败，说明数据确实包含互补信息；二是特征筛选显著影响结果，说明当前任务仍然处在高维小样本主导的阶段。

从这个角度看，本文并不追求“某次实验做到最高分”，而更关注“在固定工程框架下，各模块的作用是否可以被清晰验证”。这也是为什么消融实验比单次最优结果更值得写进论文主体，因为它更能说明方法链条的必要性。
